\chapter{Completion and coefficients}
\label{ch:completion-coefficients}

Let $k$ be a field of characteristic zero, and let $h,x$ be commuting
indeterminates. The equation $(1-hx)f=1$ has the finite approximations
\begin{equation}\label{eq:cc-geometric-finite}
 (1-hx)\sum_{j=0}^{N-1}h^jx^j=1-h^Nx^N,\qquad N\geq1.
\end{equation}

The parameter $h$ has degree zero in every complex below.
No numerical value or absolute value is assigned to it.
The remainder disappears at every fixed parameter order as $N$ increases.
To obtain an actual solution, the allowed expressions must include the
resulting series. The same choice also determines which maps extend and
whether taking cohomology preserves these limits.

\section{Polynomial coefficients and compatible truncations}
\label{sec:cc-truncations}

The partial sums in \eqref{eq:cc-geometric-finite} have unbounded degree
in $x$, although their coefficient at each power of $h$ is a polynomial.
This distinction separates two rings of expressions.
Use the formal multiplication defined in Chapter~\ref{ch:products-coefficients},
with $h$ as parameter, and put
\[
 R=k[[h]],\qquad C=k[x][[h]].
\]

An element of $R[x]$ has one finite bound on its $x$-degree.
An element of $C$ has the form $\sum_{n\geq0}h^np_n(x)$, with
$p_n\in k[x]$ and no common bound on $\deg p_n$.
For example, $\sum_{n\geq0}h^nx^{n^2}$ belongs to $C$.
At order $h^N$, multiplication in $C$ uses only
$\sum_{i=0}^N p_i(x)q_{N-i}(x)$, which is a polynomial.

The same construction applies to a $k$-vector space $V$. Write
\[
 V[[h]]=\prod_{n\geq0}h^nV
 =\left\{\sum_{n\geq0}h^nv_n:v_n\in V\right\}.
\]
The product permits an arbitrary coefficient at each index.
Scalar multiplication by $R$ uses the finite coefficient convolution.
The subspace $h^NV[[h]]$ consists of series with zero coefficients below
order $N$. Equality modulo this subspace means agreement in the first
$N$ coefficients.

Extending a coefficient formula requires control of its output precision
in terms of input precision.
A topology on a set $X$ is a collection of subsets, called open sets,
containing $X$ and the empty set.
It must be closed under arbitrary unions and finite intersections.
A neighborhood of $x\in X$ is a subset containing an open set that
contains $x$.
A set with a specified topology is a topological space.

\begin{definition}\label{def:cc-filtered-topology}
Let $P$ be a module over a commutative ring $S$.
A decreasing filtration of $P$ is a sequence of submodules
$F^0P\supseteq F^1P\supseteq\cdots$.
For $x\in P$, the coset $x+F^NP$ is the set
$\{x+v:v\in F^NP\}$.
The filtered topology has the cosets $x+F^NP$, for $N\geq0$, as
basic neighborhoods of each $x\in P$.
Explicitly, a subset $O\subseteq P$ is open when each $x\in O$ has
some $N$ with $x+F^NP\subseteq O$.
No condition $F^0P=P$ is required here.
\end{definition}

For $V[[h]]$, take $F^NV[[h]]=h^NV[[h]]$.
This filtered topology is the $h$-adic topology.

Each coset $x+F^NP$ is open, because $y+F^NP=x+F^NP$ whenever
$y\in x+F^NP$.
These open sets form a topology. Unions satisfy the condition at
each point using any member that contains it.
For two open sets containing $x$, the larger of their two cutoff
indices supplies a coset inside their intersection.
This proves the finite-intersection condition by induction.
The whole set and the empty set satisfy the defining condition as well.

A sequence $x_j$ converges to $x$ when every neighborhood of $x$
contains all sufficiently late terms.
In the filtered topology this means that, for every $N$, eventually
$x_j-x\in F^NP$.
For $V[[h]]$, it means that each finite initial list of coefficients
eventually stabilizes at the corresponding coefficients of the limit.

A map $f:X\to Y$ between topological spaces is continuous at $x$ if
each neighborhood $O$ of $f(x)$ contains $f(W)$ for some neighborhood
$W$ of $x$.
It is continuous if this holds at every point.

For modules $P,Q$ over a common commutative ring $S$, give them
filtrations $F,G$, respectively.
An arbitrary map $f:P\to Q$ is therefore continuous at $x$ exactly when
\begin{equation}\label{eq:cc-continuity-at-point}
 \text{for every }N\text{ there is }M\text{ such that}\quad
 f(x+F^MP)\subseteq f(x)+G^NQ.
\end{equation}
If $f$ is $S$-linear, the identity $f(x+v)=f(x)+f(v)$ shows that
continuity at zero is equivalent to continuity everywhere.
In that case the criterion is
$f(F^MP)\subseteq G^NQ$ for a suitable $M$ at each $N$.
Thus preserving each filtration submodule is a sufficient continuity
condition, but different input and output cutoffs are permitted.

Continuous maps preserve convergent sequences.
Indeed, for a neighborhood $O$ of $f(x)$, continuity gives a
neighborhood $W$ of $x$ with $f(W)\subseteq O$.
If $x_j\to x$, then $x_j\in W$ eventually, so $f(x_j)\in O$
eventually. This proves $f(x_j)\to f(x)$.
Composition preserves continuity by applying the neighborhood condition
successively to the two maps.

The discrete topology on a set declares every subset open.
Its basic neighborhoods are singletons, so a sequence converges exactly
when it eventually equals its limit.
For a module this is the filtered topology with $F^NP=0$ for every $N$.

For a family of topological spaces $(P_i)_{i\in I}$, their product
topology has the following basic neighborhoods of a point $(x_i)$.
Choose a finite subset $J\subseteq I$ and an open set $W_i$ containing
$x_i$ for each $i\in J$.
The resulting neighborhood consists of all $(y_i)$ with $y_i\in W_i$
for $i\in J$, with no restriction on the remaining coordinates.

Open sets in the product are unions of these basic neighborhoods.
The intersection of two basic neighborhoods combines their finite
coordinate lists and intersects the selected open sets in each factor.
The empty coordinate list gives the whole product.
Distributing finite intersections over unions proves the
finite-intersection axiom for open sets.
Thus the product construction satisfies the topology axioms.

If the factors are filtered modules, one can take
$W_i=x_i+F_i^{N_i}P_i$.
For finitely many factors a common cutoff $M\geq N_i$ suffices in
every coordinate, since the filtrations decrease.
A sequence in a product converges exactly when every coordinate
sequence converges. Necessity follows by restricting one coordinate.
For sufficiency, a basic neighborhood restricts only finitely many
coordinates, so one index beyond all their required indices works.
In particular, the $h$-adic topology on $V[[h]]$ is the product
topology with discrete coefficient spaces, because every finite set
of coefficient indices lies in one initial segment.

Addition in a filtered module is continuous: perturbing two inputs
by elements of $F^NP$ changes their sum by an element of $F^NP$.
This is the continuity criterion for the product neighborhoods on
the pair of inputs.
Negation preserves every $F^NP$, so it is continuous as well.
It follows that a finite sum of continuous maps into a filtered
module is continuous, by intersecting their finitely many required
input neighborhoods at each point.

For a subset $U\subseteq P$, its closure consists of the points
whose every neighborhood meets $U$.
Write $U+F^NP=\{u+v:u\in U,\ v\in F^NP\}$.
Since each $F^NP$ is closed under additive inverses, the closure is
\[
 \overline U=\bigcap_{N\geq0}(U+F^NP).
\]
The subset $U$ is dense in $P$ when $\overline U=P$.
Equivalently, every $x\in P$ can be approximated by an element of
$U$ modulo each $F^NP$.
These definitions also apply to the finite products above, using the
common-cutoff submodules $\prod_iF_i^MP_i$.

The polynomial module $V[h]$ consists of finite sums in $h$.
Its quotient $V[h]/h^NV[h]$ retains the coefficients from $0$ to $N-1$.
The ring $R$ acts on this quotient through coefficient truncation
$R\to k[h]/(h^N)$, followed by the finite coefficient convolution.
Reduction gives maps from level $N+1$ to level $N$.
The compatibility of all these reductions expresses completion without
an infinite summation inside $V$.

\begin{definition}\label{def:cc-inverse-limit}
An inverse system of modules over a commutative ring $S$ consists of
modules $M_N$, for $N\geq1$, and $S$-linear maps
$\rho_N:M_{N+1}\to M_N$. Its inverse limit is
\[
 \varprojlim_N M_N
 =\{(m_N)_N\in\prod_{N\geq1}M_N:
       \rho_N(m_{N+1})=m_N\text{ for every }N\}.
\]
Transitions over several levels mean composites of the displayed maps.
\end{definition}

For an $S$-module $W$, compatible $S$-linear maps $f_N:W\to M_N$
give the unique $S$-linear map $w\mapsto(f_N(w))_N$ to the inverse limit.
The module operations in the limit are coordinatewise, so this formula
proves existence, linearity, and uniqueness. This is its universal property.

\begin{proposition}\label{prop:cc-series-limit}
Coefficient reduction gives an isomorphism of $R$-modules
\[
 V[[h]]\longrightarrow\varprojlim_N V[h]/h^NV[h].
\]
For $V=k[x]$, it also identifies the rings
$C=\varprojlim_N R[x]/h^NR[x]$.
\end{proposition}

\begin{proof}
Every series gives compatible reductions. If all reductions vanish,
every coefficient vanishes, so this map is injective.

For a compatible family, take its coefficient of $h^n$ at any level
$N>n$. Compatibility makes this independent of the level.
The resulting sequence $(v_n)$ gives a series with the required reductions.
This constructs an inverse, which respects addition and scalar multiplication.
For polynomial coefficients, both rings at level $N$ consist of the
same finite lists of polynomials, with the same truncated multiplication.
The isomorphism therefore respects products as well.
\end{proof}

For decreasing submodules $F^NP\subset P$, the completion is
$\widehat P=\varprojlim_N P/F^NP$.
The module $P$ is separated if $\bigcap_N F^NP=0$, and complete if
the canonical map $P\to\widehat P$ is surjective.
Thus separatedness and completeness together mean that this map is an
isomorphism. Proposition~\ref{prop:cc-series-limit} proves both properties
for $V[[h]]$, without a dimension restriction on $V$.

Separatedness gives uniqueness of limits.
If $x_j$ converges to both $x$ and $y$, then eventually
$x_j-x,x_j-y\in F^NP$ for every fixed $N$.
Subtraction gives $x-y\in F^NP$ for every $N$, hence $x=y$.

For a submodule $U\subseteq P$, its induced topology uses the
intersections $U\cap F^NP$ as neighborhoods of zero.
This is precisely the topology obtained by intersecting neighborhoods
in $P$ with $U$.

Two continuous maps $f,g:P\to Q$ into a separated filtered module
that agree on a dense subset $U$ agree everywhere.
Fix $x\in P$ and an output cutoff $N$.
By \eqref{eq:cc-continuity-at-point}, choose one input cutoff $M$
that works for both maps at $x$.
Density gives $u\in U\cap(x+F^MP)$, so
\[
 f(x)-g(x)=(f(x)-f(u))+(g(u)-g(x))\in G^NQ.
\]
Separatedness of $Q$ now gives $f(x)=g(x)$.
The same argument applies to finite products of filtered modules,
using their common-cutoff neighborhoods.

Writing $p_n(x)=\sum_{j\geq0}a_{n,j}x^j$ displays the condition in $C$:
each fixed row $n$ of $(a_{n,j})$ has finite support in $j$.
The ring $k[[h,x]]$ instead permits an arbitrary array of coefficients.
Its coefficient of $h^nx^j$ in a product uses the finite rectangle
$0\leq a\leq n$, $0\leq b\leq j$, so multiplication is defined.
The series $\sum_{j\geq0}x^j$ belongs to $k[[h,x]]$ and does not belong
to $C$. Its row at parameter order zero is infinite.

\section{Localization and the lower bound on parameter order}
\label{sec:cc-localization}

Allowing finitely many negative powers of $h$ makes it possible to
evaluate a polynomial at $x=h^{-1}$.
The number of permitted negative orders must still be bounded for each
individual expression.

For a commutative ring $S$ and $s\in S$, define $S[1/s]$ using fractions
$a/s^m$, where $a\in S$ and $m\geq0$.
The equality rule is
\[
 \frac{a}{s^m}=\frac{b}{s^n}
 \quad\Longleftrightarrow\quad
 s^\ell(s^na-s^mb)=0\text{ for some }\ell\geq0.
\]

Addition uses a common denominator, and multiplication multiplies
numerators and denominators. Multiplying the equality equations by
further powers of $s$ shows that these operations respect representatives.
The same multiplication proves transitivity of the equality rule.
For example, equalities between denominators $s^m,s^n$ and $s^n,s^p$
give, after adding suitably multiplied equations, an equality between
denominators $s^m,s^p$, with an extra factor $s^n$.
The ring laws follow by expressing any finite calculation over a common
denominator and using the laws in $S$.

The resulting ring is the localization at $s$.
The map $S\to S[1/s]$ sends $a$ to $a/1$.
It can have a kernel when $s$ kills a nonzero element.

If $f:S\to T$ is a unital ring map and $f(s)$ is invertible, the unique
extension sends $a/s^m$ to $f(a)f(s)^{-m}$.
The equality rule proves that this is well-defined, because $f(s)^\ell$
can be cancelled. Conversely, preservation of products forces this formula.
This proves the mapping property of localization, including the case of
the zero ring, whose identity is zero.

For a nonzero series in $R$, let its order be its first nonzero exponent.
The first nonzero coefficient of a product is the product of the first
nonzero coefficients. Thus $R$ has no zero divisors.
Every nonzero element has the form $h^m v$ with $v$ a unit, by
Lemma~\ref{lem:pc-units}. Consequently
\[
 K=R[1/h]=k((h))
 =\left\{\sum_{n\geq n_0}a_nh^n:
          a_n\in k,\ n_0\in\mathbb Z\right\}
\]
is a field: the inverse of $h^mv$ is $h^{-m}v^{-1}$.

The two $K$-algebras of interest are
\[
 A=K[x],\qquad B=C[1/h]=k[x][[h]][1/h].
\]
Both contain $R[x]$, and preserving coefficients gives an inclusion
$A\hookrightarrow B$. An element of $B$ is a series
$\sum_{n\geq n_0}h^np_n(x)$ with one lower bound $n_0$ and polynomial
coefficients at every order.

For a nonzero element of $K$ or $B$, write $\operatorname{val}_h$ for
its least nonzero parameter exponent, and put
$\operatorname{val}_h(0)=+\infty$.
Regrouping by powers of $x$ gives a ring injection $B\hookrightarrow K[[x]]$.
Indeed, each $x$-coefficient is a Laurent series with the same lower
bound. Equality and multiplication in the two descriptions agree at
each pair of exponents.

\begin{proposition}\label{prop:cc-valuation-carrier}
A series $\sum_{j\geq0}b_j(h)x^j\in K[[x]]$ belongs to $B$ exactly
when, for every integer $N$, only finitely many $j$ satisfy
$\operatorname{val}_h b_j<N$.
It belongs to $A$ exactly when only finitely many $b_j$ are nonzero.
\end{proposition}

\begin{proof}
Suppose $b=h^{-m}\sum_{n\geq0}h^np_n(x)$ belongs to $B$.
Below order $N$, only finitely many polynomials $p_n$ contribute.
The union of their finite supports is finite, giving the stated condition.

Conversely, the condition at $N=0$ leaves only finitely many nonzero
coefficients with negative valuation.
Their valuations together with zero have an integer minimum $n_0$.
All coefficients therefore have order at least $n_0$.
Expand them above this common lower bound.

At a fixed order $n\geq n_0$, a nonzero coefficient can occur only for
$j$ with $\operatorname{val}_h b_j<n+1$.
There are finitely many such $j$, so this row defines a polynomial
$p_n(x)$. The expression $\sum_{n\geq n_0}h^np_n(x)$ belongs to $B$
and gives the prescribed series in $K[[x]]$.
The characterization of $A$ is precisely the finite-support definition
of a polynomial over $K$.
\end{proof}

The inclusions
\[
 A\subsetneq B\subsetneq k[[h,x]][1/h]\subsetneq K[[x]]
\]
are strict, as shown by the respective elements
\[
 \sum_{j\geq0}h^jx^j,\qquad
 \sum_{j\geq0}x^j,\qquad
 \sum_{j\geq0}h^{-j}x^j.
\]
The middle larger ring requires a common lower parameter bound but
allows infinite rows: multiplying by a power of $h$ produces precisely
an arbitrary array with nonnegative exponents in both variables.
The final series has no such lower bound.

To complete $A$, use the additive neighborhoods $h^NR[x]$, with
$N\geq0$. On $B$, use $h^NC$.
These define a lattice topology, meaning the topology specified by
powers of $h$ times the indicated $R$-submodule.
They are not ideals of the $K$-algebras: $h^NA=A$ and $h^NB=B$.
In particular, the ideal quotients $A/h^NA$ are zero and cannot produce
this completion.

\begin{proposition}\label{prop:cc-lattice-complete}
The lattice topology on $B$ is separated and complete.
Its restriction to $A$ has neighborhoods $h^NR[x]$, and $A$ is dense
in $B$. Thus $B$ is the completion of $A$ for this specified topology.
\end{proposition}

\begin{proof}
One has $A\cap h^NC=h^NR[x]$: a polynomial lies in the intersection
exactly when each of its finitely many coefficients has order at least $N$.
Also $\bigcap_Nh^NC=0$, because every nonzero series has a least exponent.

A sequence is Cauchy when, for each $N$, sufficiently late differences
belong to $h^NC$. A Cauchy sequence has a tail in one translate of $C$,
by taking $N=0$. That tail has a common lower parameter bound.
At every fixed order its polynomial coefficient eventually stabilizes.
These stable coefficients define an element of $B$, and the sequence
converges to it by the Cauchy condition.

Truncating any element of $B$ above parameter order $N-1$ gives a
finite sum of polynomials, hence an element of $A$.
The truncation error belongs to $h^NC$, proving density.

For completeness in the inverse-limit sense, truncate a representative
of each coset in $B/h^NC$ below order $N$.
This is independent of the representative, since adding an element
of $h^NC$ leaves all lower coefficients unchanged.
The truncations of compatible cosets form a Cauchy sequence.
Its limit has all the specified cosets, proving surjectivity onto
the inverse limit.
The intersection calculation proves injectivity.
Finally, the same truncation gives $A/h^NR[x]\cong B/h^NC$ as additive
groups, so their inverse limits identify this completion with that of $A$.
\end{proof}

Addition and multiplication are continuous for these topologies.
The order of a sum is at least the minimum of the orders of its terms.
For nonzero factors in $B$, the order of their product is the sum of
their orders, because the leading polynomial coefficients have nonzero
product. Addition therefore satisfies
\eqref{eq:cc-continuity-at-point} with the same input and output cutoff.

For multiplication at $(b,c)$ and output cutoff $N$, choose $M\geq0$
such that $M+\operatorname{val}_h b\geq N$,
$M+\operatorname{val}_h c\geq N$, and $2M\geq N$.
Conditions involving a zero factor impose no restriction because its
valuation is $+\infty$.
For $\delta,\epsilon\in h^MC$, the difference
\[
 (b+\delta)(c+\epsilon)-bc=b\epsilon+\delta c+\delta\epsilon
\]
belongs to $h^NC$ by these three bounds.
This proves continuity for the product neighborhoods on $B\times B$.
Intersecting with $A$ proves the same assertion there.

For a linear map $T:M\to N$, its cokernel is the quotient module
$\operatorname{coker}T=N/T(M)$.
It measures the output classes modulo outputs attained by $T$.

\begin{proposition}\label{prop:cc-geometric}
Multiplication by $1-hx$ is an injective $K$-linear map $A\to A$ with
cokernel $K$. On $B$ it is invertible, with inverse multiplication by
\[
 g=\sum_{n\geq0}h^nx^n\in C.
\]
Evaluation $f\mapsto f(h^{-1})$ on $A$ has no extension to a unital
$K$-algebra map $B\to K$.
\end{proposition}

\begin{proof}
If $f\in A$ is nonzero of degree $d$, the leading term of $(1-hx)f$
has degree $d+1$ and nonzero coefficient. This proves injectivity.
Evaluation at $h^{-1}$ is surjective because it fixes the constants $K$.
For $a=h^{-1}$, the identity
\[
 x^j-a^j=(x-a)\sum_{i=0}^{j-1}x^{j-1-i}a^i\qquad(j\geq1)
\]
shows that $f(x)-f(a)$ is a multiple of $x-a$.
Thus its kernel is $(x-h^{-1})A=(1-hx)A$, since the two generators
differ by the unit $-h$. Evaluation therefore identifies the cokernel
$A/(1-hx)A$ with $K$.

The series $g$ belongs to $C$ by its polynomial coefficients.
Equation~\eqref{eq:cc-geometric-finite} gives $(1-hx)g=1$
coefficientwise. Commutativity gives a two-sided inverse in $B$.
An extension of evaluation would send this unit to $1-hh^{-1}=0$.
Applying the extension to $(1-hx)g=1$ would give $0=1$ in $K$,
which is impossible.
\end{proof}

Give $K$ the lattice neighborhoods $h^NR$, for $N\geq0$.
Evaluation is already discontinuous from $A$ to $K$ with these topologies:
$h^nx^n\to0$, while every evaluation equals $1$.
The nonextension assertion is stronger, because it excludes even a
discontinuous unital algebra extension.

\section{A differential equation and its inverse}
\label{sec:cc-differential}

The preceding multiplication map loses a cokernel after completion.
The differential equation $(\partial_x-h)f=0$ instead gains solutions.
Define polynomial differentiation by $\partial_x(x^j)=jx^{j-1}$ for
$j\geq1$ and $\partial_x(1)=0$.
Apply it coefficientwise in $h$ on $C$ and $B$.
It preserves $h^NC$, so
$D=\partial_x-h:B\to B$ is continuous and $K$-linear.
It preserves $A$ as well.

Let $M$ be a module over a commutative ring $S$, and let
$T:M\to M$ be $S$-linear.
Write $H^q(M,T)$ for the degree-$q$ cohomology of the two-term complex
$[M\xrightarrow{T}M]$, concentrated in degrees zero and one.
Thus
\[
 H^0(M,T)=\ker T,\qquad H^1(M,T)=M/T(M).
\]
The second group is $\operatorname{coker}T$, and all other cohomology
groups vanish.
Proposition~\ref{prop:cc-geometric} consequently gives $H^1=K$ on $A$
and an acyclic complex on $B$ for $T=1-hx$.

\begin{theorem}\label{thm:cc-differential}
For $D=\partial_x-h$, one has
\[
 H^0(A,D)=H^1(A,D)=0,\qquad
 H^0(B,D)=K E,\qquad H^1(B,D)=0,
\]
where $E=\exp(hx)$ denotes the formal series
$\sum_{n\geq0}h^nx^n/n!\in C$.
\end{theorem}

\begin{proof}
For $q\in A$ of degree at most $d$, define
\[
 S(q)=-\sum_{j=0}^d h^{-j-1}\partial_x^jq.
\]
The sum is finite, and cancellation gives $DS(q)=q$.
Increasing $d$ does not change it, because higher derivatives vanish.
For a nonzero polynomial $f$, the term $-hf$ has its original degree,
whereas $\partial_xf$ has smaller degree. Thus $Df\ne0$.
Surjectivity and injectivity prove that $S$ is the inverse of $D$ on $A$.

Given $q=\sum_{n\geq0}h^nq_n(x)\in C$, seek
$f=\sum_{n\geq0}h^nf_n(x)$ with $Df=q$.
The coefficient equations are
\[
 f_n'=q_n+f_{n-1},\qquad n\geq0,\qquad f_{-1}=0.
\]
Every polynomial has an antiderivative over $k$, since every positive
integer is invertible. At each step choose the antiderivative with zero
constant term. This constructs $f\in C$ and proves surjectivity there.

For $q\in B$, multiply by a power of $h$ to place it in $C$, solve,
and divide by the same power. This proves surjectivity on $B$.

Differentiating the coefficients of $E$ gives $\partial_xE=hE$.
The series $\exp(-hx)$ is its inverse: its product with $E$ has
coefficient $x^n\sum_{i=0}^n(-1)^i/(i!(n-i)!)$ at order $h^n$.
This coefficient is zero for $n>0$ by the finite binomial expansion
of $(1-1)^n$, and it is one for $n=0$.

The polynomial product rule extends at each parameter coefficient.
It therefore gives
\[
 D(Ef)=E\partial_xf\qquad(f\in B).
\]
A polynomial over $k$ has zero derivative exactly when it is constant.
Applied to every parameter coefficient, this proves $\ker\partial_x=K$
on $B$. Multiplication by the unit $E$ identifies that kernel with
$\ker D$, giving $K E$.
\end{proof}

The inverse $S:A\to A$ is discontinuous for the induced lattice topology.
Indeed, $q_m=h^mx^m\to0$, but
\[
 S(q_m)=-\sum_{j=0}^m\frac{m!}{(m-j)!}h^{m-j-1}x^{m-j}
\]
has order $-1$ for every $m$. Its constant $x$-coefficient is
$-m!h^{-1}$, which is nonzero in characteristic zero.
Thus an algebraic inverse on polynomials need not extend continuously.

The new kernel element also comes from finite approximations.
With $E_N=\sum_{j=0}^Nh^jx^j/j!$, direct cancellation gives
$DE_N=-h^{N+1}x^N/N!$. Hence $E_N\to E$ and $DE_N\to0$.

Characteristic zero is essential here. In characteristic $p>0$,
$\partial_x^p=0$ on $k[x]$ and coefficientwise on $B$.
Indeed, a nonzero $p$-fold monomial derivative would have a coefficient
containing $p$ consecutive integer factors, one divisible by $p$.
Then $-\sum_{j=0}^{p-1}h^{-j-1}\partial_x^j$ is a two-sided inverse
of $D$ on $B$, by the same finite cancellation.
The exponential formula in the theorem has undefined denominators in
that characteristic and cannot produce the asserted kernel.

\section{Solving one coefficient at a time}
\label{sec:cc-recursion}

The antiderivative construction solves each coefficient equation after
the preceding coefficients have been chosen.
A general formal operator has the same form whenever its parameter
powers are nonnegative.
Let $S$ be a commutative ring, let $V,W$ be $S$-modules, and let
$T_i:V\to W$ be $S$-linear maps for $i\geq0$.
Their formal operator is the map
\[
 T(h):V[[h]]\longrightarrow W[[h]],\qquad
 [h^n]T(h)v=\sum_{i=0}^nT_i(v_{n-i}).
\]

The series modules and their filtrations have the same coefficient
definition as above, now over $S$.
Each displayed sum is finite. Distributing the finite convolutions
proves that $T(h)$ is $S[[h]]$-linear and preserves the filtration.

\begin{theorem}\label{thm:cc-recursive-inverse}
If $T_0$ is bijective, then $T(h)$ has a unique inverse.
The inverse is $S[[h]]$-linear and preserves every submodule
$h^NW[[h]]$, with its image in $h^NV[[h]]$.
In particular, it is continuous. No freeness or finite-generation
assumption on $V,W$ is needed.
\end{theorem}

\begin{proof}
For $w=\sum_{n\geq0}h^nw_n$, the equation $T(h)v=w$ forces
\begin{equation}\label{eq:cc-recursive-inverse}
 v_0=T_0^{-1}w_0,\qquad
 v_n=T_0^{-1}\left(w_n-\sum_{i=1}^nT_i(v_{n-i})\right)
 \quad(n\geq1).
\end{equation}
Conversely, these equations construct one coefficient $v_n$ at each
step, and substitution proves $T(h)v=w$.
Every solution must have these coefficients, proving uniqueness.
Since $T(h)$ is $S[[h]]$-linear, uniqueness makes its inverse linear
over the same ring.
If $w_n=0$ for $n<N$, induction in
\eqref{eq:cc-recursive-inverse} gives $v_n=0$ for $n<N$.
This proves the stated filtration bound and continuity.
\end{proof}

For an $S$-linear map $L:V\to V$, this theorem gives
\[
 (1-hL)^{-1}=\sum_{j\geq0}h^jL^j
 \quad\text{on }V[[h]].
\]
At order $n$ only $j\leq n$ contributes.
On the polynomial module $V[h]$, the operator $1-hL$ is bijective
exactly when $L$ is locally nilpotent: for each $v\in V$, some
power $L^m(v)$ is zero.
Under that condition the inverse series terminates on every polynomial
input, since it has only finitely many input coefficients.
Conversely, a polynomial solution to $(1-hL)f=v$ for a constant input
$v$ must have coefficients $f_j=L^jv$.
Its finite support forces these powers to vanish eventually.

Under the hypotheses of Theorem~\ref{thm:cc-recursive-inverse},
write $U=T(h)^{-1}$.
Define $V((h))=V[[h]][1/h]$ as coefficient series with one finite
lower parameter bound, and define $W((h))$ similarly.
The localized maps have opposite domains:
\[
 \begin{aligned}
 T_{\mathrm{loc}}&:V((h))\longrightarrow W((h)),\\
 T_{\mathrm{loc}}(h^{-m}v)&=h^{-m}T(h)v,\\
 U_{\mathrm{loc}}&:W((h))\longrightarrow V((h)),\\
 U_{\mathrm{loc}}(h^{-m}w)&=h^{-m}U(w),
 \end{aligned}
\]
where $m\geq0$, $v\in V[[h]]$, and $w\in W[[h]]$.

If two representations of a Laurent input agree, multiplication by
a common power of $h$ makes them equal as coefficient series.
Both $T(h)$ and $U$ commute with this multiplication, so the
displayed values are independent of the representations.
The identities $UT(h)=1$ and $T(h)U=1$ give both localized inverse
identities after clearing that same common power.
These maps preserve the respective lattice neighborhoods
$h^NV[[h]]$ and $h^NW[[h]]$.
For the polynomial inverse $S$ in Theorem~\ref{thm:cc-differential},
the negative parameter orders instead grow without bound with the input degree.

The sign of the parameter powers is essential to the recursive theorem.
For example, on $k((h))^2$ the matrix
\[
 \begin{pmatrix}1&h\\h^{-1}&1\end{pmatrix}
\]
has identity constant coefficient and annihilates the nonzero vector
$(-h,1)$. It is not a formal operator with only nonnegative powers.
Thus an invertible coefficient at order zero alone gives no inverse
when negative operator powers are also permitted.

\begin{proposition}\label{prop:cc-recursive-right-inverse}
Suppose $T_0$ has a specified $S$-linear right inverse $J:W\to V$,
so $T_0J=1_W$. Then $T(h)$ is surjective and has a continuous
$S[[h]]$-linear right inverse.
For any $w=\sum h^nw_n$ and any sequence $a_n\in\ker T_0$, all
solutions are given uniquely by the recursion
\[
 v_n=J\left(w_n-\sum_{i=1}^nT_i(v_{n-i})\right)+a_n,
 \qquad n\geq0.
\]
The sum is empty when $n=0$.
\end{proposition}

\begin{proof}
Applying $T_0$ proves the coefficient equation at order $n$.
Induction therefore constructs a solution for each sequence $(a_n)$.
For an arbitrary solution, subtracting the displayed $J$-term defines
the unique possible $a_n$. The equation at order $n$ gives $T_0a_n=0$.
This proves the parameterization.

Choose $a_n=0$ for every $n$ and denote the resulting solution map by
$U:W[[h]]\to V[[h]]$.
The recursion makes $U$ additive and $S$-linear.
It commutes with multiplication by $h$, since shifting the input
coefficients shifts the recursion, with zero constant coefficient.
It preserves every $h^N$ submodule by induction.
Hence it is continuous and commutes with multiplication by a general
series in $S[[h]]$, by taking the limits of its polynomial truncations.
The construction gives $T(h)U=1$.
\end{proof}

For $T_0=\partial_x$ on $k[x]$,
$J(x^j)=x^{j+1}/(j+1)$ is a right inverse.
Here $a_n\in k$ records the integration constant at order $n$
in the equation $Df=q$.

\section{Length restrictions and continuous operations}
\label{sec:cc-length}

The polynomial restriction can also occur in a word length or another
nonnegative index. Let $V_\ell$ be a $k$-vector space for each
$\ell\geq0$. The direct sum $V=\bigoplus_{\ell\geq0}V_\ell$ permits
only finitely many nonzero components in each element.
Its parameter series form the middle term in
\[
 \bigoplus_{\ell\geq0}V_\ell[[h]]
 \ \subseteq\
 \left(\bigoplus_{\ell\geq0}V_\ell\right)[[h]]
 \ \subseteq\
 \prod_{\ell\geq0}V_\ell[[h]].
\]

The first space has one finite support over all parameter orders.
The middle space has finite support at each order.
The last space permits arbitrary support.
In the middle space, an equivalent condition on components
$v_\ell(h)\in V_\ell[[h]]$ is
\[
 \text{for each }N\geq1,\quad
 \#\{\ell:v_\ell(h)\notin h^NV_\ell[[h]]\}<\infty.
\]
Indeed, the union of the finite supports in rows $0,\ldots,N-1$ is
finite. Conversely, this condition at $N=n+1$ makes row $n$ finite.

For $V_\ell=k e_\ell$, the family with components $h^\ell e_\ell$
belongs to the middle space and not the first.
The family with components $e_\ell$ belongs to the last space and
not the middle. Both inclusions are therefore strict in this example.

A finite multilinear operation respects parameter order without any
uniform bound on word length.
Let $r\geq1$ be an integer and let
$m:V_1\times\cdots\times V_r\to W$ be a $k$-multilinear map.
Here multilinear means linear in each separate input.

\begin{proposition}\label{prop:cc-multilinear-extension}
There is a unique continuous $R$-multilinear map
\[
 \widehat m:V_1[[h]]\times\cdots\times V_r[[h]]\to W[[h]]
\]
that agrees with $m$ on constant inputs.
Continuity uses the product of the $h$-adic input topologies.
The extension is
\[
 [h^n]\widehat m(v_1,\ldots,v_r)
 =\sum_{n_1+\cdots+n_r=n}
       m(v_{1,n_1},\ldots,v_{r,n_r}),
 \qquad n_i\geq0.
\]
\end{proposition}

\begin{proof}
For fixed $n$, there are only finitely many tuples of nonnegative
integers with sum $n$. Thus the coefficient belongs to $W$.
The linearity identities follow by distributing these finite sums.
In particular, multiplication of an input by an element of $R$ gives
the same scalar convolution in the output.
If one input changes by an element of $h^NV_i[[h]]$, the output
changes by an element of $h^NW[[h]]$.
Replacing the inputs one at a time proves the same conclusion when
all inputs change in this way, and proves joint continuity.

An $R$-multilinear extension has the specified values on polynomial
inputs, by expansion into constants and powers of $h$.
Their simultaneous truncations converge to arbitrary series inputs.
Continuity and separatedness of $W[[h]]$ therefore force the displayed
values on all inputs, proving uniqueness.
\end{proof}

If $W$ is itself a direct sum over lengths, each value of $m$ has
finite support by its specified codomain.
Every coefficient of $\widehat m$ is a finite sum of such values,
so this support condition persists.

Finite multilinear identities formed from finite sums and compositions
also persist. At any fixed parameter order, expansion gives a finite
sum of the original identities.

For a graded space $V=\bigoplus_qV^q$, completion here means
$(V[[h]])^q=\prod_{n\geq0}h^nV^q$, with a direct sum over total degrees
$q$ as the underlying graded space.
For such spaces, each operation has its specified degree and uses
the original grading signs. The degree-zero parameter introduces no
additional signs.

For an ungraded vector space $W$, a bound on parameter order suffices
to sum infinitely many operations.
Suppose $X$ is a set and maps $f_j:X\to W[[h]]$, for $j\geq0$,
satisfy $f_j(X)\subseteq h^{a_j}W[[h]]$.
Here the integers $a_j\geq0$ tend to infinity, meaning that for
each $N$ only finitely many $j$ satisfy $a_j<N$.
At every coefficient only finitely many summands contribute, so their
sum defines a map $f=\sum_jf_j:X\to W[[h]]$.

Give $W[[h]]$ its $h$-adic topology and suppose $X$ is topological.
If every $f_j$ is continuous, so is $f$.
For an output cutoff $N$, put $J_N=\{j:a_j<N\}$.
Then $f(x)-\sum_{j\in J_N}f_j(x)$ belongs to $h^NW[[h]]$ for every $x$.
Continuity of the finite sum gives a neighborhood $O$ of $x$ on
which its difference from its value at $x$ lies in $h^NW[[h]]$.
For $y\in O$, the two discarded tails lie in this same submodule,
so $f(y)-f(x)\in h^NW[[h]]$, proving continuity.

For a graded target $W=\bigoplus_qW^q$, parameter order must be
combined with finite support in total degree. Put
\[
 G_W=\bigoplus_{q\in\mathbb Z}W^q[[h]],\qquad
 F^NG_W=\bigoplus_{q\in\mathbb Z}h^NW^q[[h]],
\]
and give $G_W$ this filtered topology.
Its restriction to $W^q[[h]]$ is the usual $h$-adic topology.
For $w\in G_W$, its degree support is the finite set of $q$ for
which its component in $W^q[[h]]$ is nonzero.
The parameter $h$ still has degree zero.

Suppose maps $f_j:X\to G_W$ satisfy $f_j(X)\subseteq F^{a_j}G_W$,
where $a_j\geq0$ and $a_j\to\infty$.
Assume that, for each $x$, there is a finite set $D_x\subset\mathbb Z$
containing the degree support of every $f_j(x)$.
No uniform bound on $D_x$ as $x$ varies is required.
The coefficient sum has support contained in $D_x$, so it belongs
to $G_W$.

If the maps are continuous, the same finite-sum argument modulo
$F^NG_W$ proves their sum continuous.
The finite degree condition ensures membership in the target, while
the uniform parameter-order bound supplies the continuity estimate.

In particular, let $X=\bigoplus_pX^p$ be a graded vector space with
a specified topology, and let every $f_j:X\to G_W$ be continuous
and $k$-linear of one fixed degree $d\in\mathbb Z$.
This degree condition means $f_j(X^p)\subseteq W^{p+d}[[h]]$.
Write $x=\sum_{p\in E_x}x_p$ for its finitely many nonzero homogeneous
components, with $x_p\in X^p$.
Linearity gives support contained in the finite set
$E_x+d=\{p+d:p\in E_x\}$ for every $f_j(x)$.
Thus the uniform parameter-order hypothesis defines a continuous
sum of degree $d$ with values in $G_W$.
It is linear because addition and scalar multiplication commute with
each finite coefficient sum.

The parameter-order condition alone does not ensure that the sum belongs to $G_W$.
Take $X=k$ with the discrete topology in degree zero and
$W=\bigoplus_{j\geq0}ke_j$ with $|e_j|=j$.
The maps $f_j(\lambda)=\lambda h^je_j$ are continuous and have
parameter order $j$.
Their coefficient sum at $1$ is $\sum_{j\geq0}h^je_j$.
Its component in every nonnegative total degree is nonzero, so it
does not belong to $G_W$.
It belongs to the coefficient-series module obtained by forgetting
the grading, which gives a larger module.

The unrestricted length product can destroy continuity even for a
linear operation. Define
\[
 b:\bigoplus_{\ell\geq0}k e_\ell\to k f,\qquad b(e_\ell)=f,
\]
where $f$ is a nonzero basis vector.
Give each coordinate space the discrete topology and the full product
the topology of agreement in each fixed finite set of coordinates.
Then $e_\ell\to0$ in that product, while $b(e_\ell)=f$ for every
$\ell$. Hence $b$ has no continuous extension from the product to
the discrete space $k f$.

Its parameter extension is defined, and satisfies
\[
 \widehat b\left(\sum_{\ell\geq0}h^\ell e_\ell\right)
       =\sum_{\ell\geq0}h^\ell f\in(k f)[[h]].
\]
The target completion and the finite support at each input order make
every coefficient a finite calculation.

\section{The algebraic tensor product and its completion}
\label{sec:cc-tensor}

The series $\sum_{n\geq0}x^ny^n$ is a limit of finite sums of products
in separate variables. An algebraic tensor product requires one finite
sum for the entire expression.
To compare these requirements, put $M=k[[x]]$ and $N=k[[y]]$.
Here $x,y$ are independent commuting variables of degree zero.

The bilinear map $(f,g)\mapsto f(x)g(y)$ induces
\[
 \alpha:M\otimes_kN\longrightarrow k[[x,y]],\qquad
       f\otimes g\longmapsto f(x)g(y).
\]
Its coefficient at $x^iy^j$ is $([x^i]f)([y^j]g)$.
The tensor product is the algebraic one from
Chapter~\ref{chap:relations-homotopies}: each element is a finite sum
of elementary tensors.

\begin{proposition}\label{prop:cc-algebraic-tensor}
The map $\alpha$ is injective.
Its image consists exactly of the coefficient arrays whose columns
span a finite-dimensional $k$-subspace of $k[[x]]$.
In particular, $\sum_{n\geq0}x^ny^n$ does not belong to its image.
\end{proposition}

\begin{proof}
Write a tensor as $\sum_{s=1}^r f_s\otimes g_s$.
Replacing the $f_s$ by a basis of their finite span and collecting
the second factors makes the $f_s$ linearly independent.
If its image is zero, its coefficient at $y^j$ gives
$\sum_s([y^j]g_s)f_s=0$ in $k[[x]]$ for every $j$.
Independence forces every coefficient of every $g_s$ to vanish.
Thus the tensor is zero, proving injectivity.

For any image element, its $j$th column belongs to the span of the
finitely many series $f_s$.
Conversely, choose a basis $f_1,\ldots,f_r$ of the column span of
an array. Write column $j$ as $\sum_s c_{s,j}f_s$ and define
$g_s(y)=\sum_{j\geq0}c_{s,j}y^j$.
Then $\sum_s f_s\otimes g_s$ maps to the array.

The columns of $\sum_nx^ny^n$ are $1,x,x^2,\ldots$, which are
linearly independent. Its column span is infinite-dimensional.
\end{proof}

To admit such limits, specify which tensor coefficients must stabilize.
For decreasing subspaces $F^iM\subset M$, $G^jN\subset N$ with
$F^0M=M$, $G^0N=N$, define
\[
 T^n=\sum_{\substack{i,j\geq0\\i+j=n}}
       \operatorname{im}(F^iM\otimes_kG^jN\to M\otimes_kN),
 \qquad
 M\widehat\otimes_kN=\varprojlim_{n\geq1}(M\otimes_kN)/T^n.
\]
The image maps come from the two subspace inclusions.
The subspaces $T^n$ decrease: in a summand with $i+j=n+1$, decrease
a positive index by one and use the corresponding input inclusion.
This specifies both the scalar base and the filtration of the
completed tensor product.
For the present $M,N$, take $F^iM=x^iM$ and $G^jN=y^jN$.

For $n\geq1$, the rectangular subspaces
\[
 U_n=x^nM\otimes_kN+M\otimes_ky^nN
 \quad\subset M\otimes_kN
\]
use the images of those tensor maps as well.
They satisfy $T^{2n-1}\subset U_n\subset T^n$.
For the first containment, if $i+j=2n-1$, either $i\geq n$ or
$j\geq n$. The second follows from the terms $(i,j)=(n,0),(0,n)$.

Put $Q=M\otimes_kN$.
The containments give the quotient maps
\[
 a_n:Q/T^{2n-1}\longrightarrow Q/U_n,\qquad
 b_n:Q/U_n\longrightarrow Q/T^n,
\]
each sending the class of $q\in Q$ to the class of the same element.
These maps are well-defined because a change by the source subspace
is also a change by the target subspace.

Define maps between the two inverse limits by
\[
 \begin{aligned}
 \Phi&:\varprojlim_mQ/T^m\longrightarrow\varprojlim_nQ/U_n,\\
 \Phi((\xi_m)_m)_n&=a_n(\xi_{2n-1}),\\
 \Psi&:\varprojlim_nQ/U_n\longrightarrow\varprojlim_mQ/T^m,\\
 \Psi((\eta_n)_n)_m&=b_m(\eta_m).
 \end{aligned}
\]
For $\Phi$, reducing coordinate $n+1$ to $n$ amounts to reducing
$\xi_{2n+1}$ modulo $U_n$.
Compatibility of $\xi$ makes this the same as reducing $\xi_{2n-1}$
modulo $U_n$, giving its required coordinate $n$.
For $\Psi$, reduction from $m+1$ to $m$ gives the class of $\eta_{m+1}$
modulo $T^m$, which equals that of $\eta_m$ by compatibility.
Thus both formulas produce compatible families.

The coordinate $(\Psi\Phi(\xi))_n$ is the reduction of $\xi_{2n-1}$
modulo $T^n$, hence equals $\xi_n$.
The coordinate $(\Phi\Psi(\eta))_n$ is the reduction of $\eta_{2n-1}$
first modulo $T^{2n-1}$ and then modulo $U_n$.
This is the direct reduction from $Q/U_{2n-1}$ to $Q/U_n$, and
compatibility gives $\eta_n$.
Both composites are identities, so the two inverse limits are isomorphic.

The coordinate projection at index $n$ sends a limit element to its
component at that index.
Give the limits the filtrations formed by the kernels of their
coordinate projections.
Output cutoff $n$ for $\Phi$ uses input cutoff $2n-1$, while output
cutoff $m$ for $\Psi$ uses input cutoff $m$, proving both maps continuous.

Write $k\{v_1,\ldots,v_r\}$ for the span of the listed vectors.
The decompositions
$M=k\{1,x,\ldots,x^{n-1}\}\oplus x^nM$ and
$N=k\{1,y,\ldots,y^{n-1}\}\oplus y^nN$
identify $(M\otimes_kN)/U_n$ with
$(M/x^nM)\otimes_k(N/y^nN)$.
Each quotient has basis $x^iy^j$ for $0\leq i,j<n$.
Compatible rectangles therefore give every coefficient array uniquely.

It follows that
\begin{equation}\label{eq:cc-completed-tensor}
 k[[x]]\widehat\otimes_k k[[y]]
 \cong\varprojlim_n k[x,y]/(x^n,y^n)
 \cong k[[x,y]].
\end{equation}
The maps take every elementary tensor to its coefficient product.
Finite rectangular polynomials are dense on the right and come from
the algebraic tensor product, so the completion map has dense image.
The topology also equals the topology of powers of $(x,y)$, since
\[
 (x,y)^{2n-1}\subseteq(x^n,y^n)\subseteq(x,y)^n.
\]
These ideal containments follow from the same exponent calculation.

The completed tensor has a useful mapping property with uniform input
precision. Let $P$ be a $k$-vector space, separated and complete for
decreasing subspaces $F^rP$, $r\geq1$.
Let $\beta:M\times N\to P$ be bilinear.
Suppose that for every $r$ there are $a,b\geq1$ such that
\[
 \beta(x^aM,N)\subset F^rP,
 \qquad \beta(M,y^bN)\subset F^rP.
\]
The bounds are independent of the other input.
Then $\beta$ induces a unique continuous linear map from the completed
tensor in \eqref{eq:cc-completed-tensor} to $P$.

To prove this, reduce the output modulo $F^rP$.
The displayed bounds make the bilinear map factor through
$(M/x^aM)\times(N/y^bN)$ and hence through its tensor product.
Any two sufficient pairs of bounds agree after reduction from a
larger rectangle. The induced maps from the completed tensor to
$P/F^rP$ are therefore well-defined and compatible as $r$ varies.
They give a map to $\varprojlim_rP/F^rP=P$.

Each reduction uses a finite rectangle, proving continuity.
Agreement on the dense algebraic tensor and separatedness of $P$
prove uniqueness.

Continuity can require different input and output cutoffs.
Differentiation on $k[[x]]$ is continuous because
$\partial_x(x^{n+1}k[[x]])\subset x^nk[[x]]$.
It does not define an endomorphism of $k[[x]]/(x^n)$ for $n\geq1$
in characteristic zero: the representative $x^n$ of zero has derivative
$nx^{n-1}$, which is nonzero in that quotient.
The correct finite map takes input modulo $x^{n+1}$ and output modulo
$x^n$.

\section{Scalar extension before and after a limit}
\label{sec:cc-scalar-limit}

Extending coefficients at every finite order supplies all coefficients
in the larger field. An algebraic scalar extension of a whole series
can still impose a finite-dimensional restriction on its coefficients.
Let $z$ be an indeterminate independent of $h$.
The rational function field $L=k(z)$ consists of fractions $a(z)/b(z)$
of polynomials, with $b\ne0$, identified by cross multiplication.
The polynomial ring has no zero divisors, so this equality is transitive.
The common-denominator operations respect equality and give a field:
a nonzero fraction $a/b$ has inverse $b/a$.

The coefficient extension map is
\[
 \gamma:L\otimes_k k[[h]]\longrightarrow L[[h]],\qquad
 a\otimes\sum_{n\geq0}c_nh^n\longmapsto\sum_{n\geq0}ac_nh^n.
\]
At each finite order the same formula gives the isomorphism
\[
 L\otimes_k\bigl(k[h]/h^Nk[h]\bigr)
       \cong L[h]/h^NL[h].
\]
Its inverse expands the finite list of coefficients as
$\sum_{n=0}^{N-1}a_n\otimes h^n$.
Proposition~\ref{prop:cc-series-limit} identifies the limit of these
finite targets with $L[[h]]$.

\begin{proposition}\label{prop:cc-scalar-extension}
The map $\gamma$ is injective. Its image consists exactly of series
whose coefficients belong to a fixed finite-dimensional $k$-subspace
of $L$. In particular, $\sum_{n\geq0}z^nh^n$ is not in its image.
\end{proposition}

\begin{proof}
Write an input as $\sum_{j=1}^r a_j\otimes f_j(h)$.
Choose a basis of the span of the $a_j$ and collect the second
factors, so that the $a_j$ are linearly independent.
If the image is zero, its coefficient at $h^n$ gives
$\sum_j([h^n]f_j)a_j=0$ for every $n$.
Independence forces all $f_j$ to vanish, proving injectivity.

Every coefficient of the image belongs to the span of the $a_j$.
Conversely, choose a basis $a_1,\ldots,a_r$ of a finite-dimensional
subspace containing all coefficients of a given series.
The coordinates of its successive coefficients define series
$f_1,\ldots,f_r\in k[[h]]$.
Their finite tensor sum maps to the given series.
Finally, $1,z,z^2,\ldots$ are linearly independent over $k$:
a finite relation would be a zero polynomial in the indeterminate $z$.
Their infinite span excludes $\sum_nz^nh^n$.
\end{proof}

To preserve a short exact sequence under scalar extension, it suffices
to express its middle space as a direct sum of its kernel and
another subspace. For unrestricted dimensions, constructing this
decomposition uses a choice principle.

Order a collection $\mathcal P$ of subsets of a fixed set by inclusion.
A chain in $\mathcal P$ is a subcollection in which any two members
are comparable: one is contained in the other.
An upper bound for a chain is a member of $\mathcal P$ containing
every set in the chain.
A member is maximal if no strictly larger member belongs to
$\mathcal P$.
We assume the following maximality form of the axiom of choice:
every nonempty such $\mathcal P$ in which every chain has an upper
bound has a maximal member.

For a subspace $U\subseteq V$, a complement is a subspace $L$ such
that $V=U\oplus L$, meaning that every $v\in V$ has a unique
expression $v=u+\ell$ with $u\in U$ and $\ell\in L$.

\begin{lemma}\label{lem:cc-basis-extension}
Under the stated maximality principle, every independent subset of a
$k$-vector space extends to a basis.
Every subspace has a complement.
These choices can be made simultaneously for a family indexed by a set.
\end{lemma}

\begin{proof}
Let $I$ be an independent subset of a vector space $V$.
Take $\mathcal P$ to be the collection of independent subsets of
$V$ containing $I$, ordered by inclusion.
It is nonempty because $I\in\mathcal P$.
For a nonempty chain, its union contains $I$ and is independent.
A finite relation with at least one term uses vectors from finitely
many chain members.

Repeated comparison selects one of these finitely many members that
contains all the others.
Independence in that member forces every coefficient to vanish.
The union is therefore an upper bound in $\mathcal P$.
The empty chain has upper bound $I$.
The assumed principle gives a maximal member $B$.

If $v\notin\operatorname{span}B$, a relation
$av+\sum_{j=1}^r a_jb_j=0$, with distinct $b_j\in B$, forces $a=0$.
Otherwise division by $a$ would place $v$ in the span of $B$.
Independence of $B$ then forces every $a_j=0$.

Thus $B\cup\{v\}$ would be a strictly larger independent set,
contradicting maximality.
Every vector consequently belongs to the span of $B$, so $B$ is a basis.
Taking $I$ empty proves basis existence, including the zero space.

For a subspace $U\subseteq V$, first choose a basis $I$ of $U$,
then extend it to a basis $B$ of $V$.
Let $L$ be the span of $B\setminus I$.
Splitting each finite basis expansion gives $V=U+L$.

An element of $U\cap L$ would give a finite relation among the
disjoint subsets $I$ and $B\setminus I$ of the basis $B$.
Independence forces it to be zero.
If $u+\ell=u'+\ell'$, then $u-u'=\ell'-\ell\in U\cap L$, so
both differences vanish.
This proves the required unique decomposition.

For simultaneous choices, let $(V_j,I_j)_{j\in J}$ be a family of
spaces with specified independent subsets, where $J$ is a set.
Use subsets of the set of labeled pairs $(j,v)$, with $j\in J$
and $v\in V_j$.
Require each subset to contain all $(j,v)$ with $v\in I_j$ and
require its vectors with each fixed label $j$ to be independent in $V_j$.
The set of all prescribed labeled pairs satisfies these conditions,
so the collection is nonempty.

Unions of chains satisfy these conditions by the same finite-relation
argument. A maximal subset spans every $V_j$, since a missing vector
in any one space could otherwise be adjoined with its label.
This gives bases extending all $I_j$ simultaneously.

For a family of inclusions $U_j\subseteq V_j$, first apply this
construction to the spaces $U_j$ with empty independent subsets.
Then extend the resulting bases simultaneously in the spaces $V_j$.
Taking the spans of the complementary basis subsets gives all the
required complements, by the preceding argument for a single inclusion.
\end{proof}

The extension $k\subseteq k(z)$ is flat, in the sense of preserving short exact
sequences under tensoring, as defined in
Section~\ref{sec:rel-coefficient-change}.
Consider a short exact sequence of vector spaces
$0\to U\xrightarrow{j}V\xrightarrow{q}W\to0$.
Lemma~\ref{lem:cc-basis-extension} gives a complement $V_0$ to $j(U)$
in $V$.
The restriction $q|_{V_0}:V_0\to W$ is injective because
\[
 V_0\cap\ker q=V_0\cap j(U)=0.
\]
It is surjective because, for any $v=j(u)+\ell$, one has $q(v)=q(\ell)$,
and $q$ itself is surjective.

The inverse of $q|_{V_0}$ is linear and defines a section $s:W\to V$
with $qs=1_W$.
The map $(u,w)\mapsto j(u)+s(w)$ is an isomorphism $U\oplus W\to V$:
the direct-sum decomposition proves both existence and uniqueness of
these coordinates.
Tensoring preserves this decomposition, its inclusion, and its
projection. Hence the resulting sequence is again short exact.

Thus flatness does not make scalar extension commute with the inverse
limit above. Since $k[[h]]$ is also the product of its coefficient
spaces $k$, the same map gives a failure to commute with an infinite
product.

A finite free tensor factor has a different result.
For a commutative ring $S$, an $S$-module $P$ is finite free when
it has a finite basis $e_1,\ldots,e_d$, including $d=0$.
The basis expresses every element uniquely as $\sum_i a_i e_i$.

\begin{theorem}\label{thm:cc-finite-free-limit}
For a finite free $S$-module $P$ and any inverse system $(M_N,\rho_N)$
of $S$-modules, the map
\[
 \lambda:P\otimes_S\varprojlim_N M_N
       \longrightarrow\varprojlim_N(P\otimes_S M_N),\qquad
 p\otimes(m_N)_N\longmapsto(p\otimes m_N)_N
\]
is an isomorphism. The target transitions are $1_P\otimes\rho_N$.
No surjectivity condition on $\rho_N$ is required.
\end{theorem}

\begin{proof}
Scalar balancing and additivity make $\lambda$ well-defined.
Using the fixed basis of $P$, write a compatible target element at
level $N$ uniquely as $\sum_{i=1}^d e_i\otimes m_{i,N}$.
Compatibility of these expressions gives
$\rho_N(m_{i,N+1})=m_{i,N}$ for every $i$.
Hence each $(m_{i,N})_N$ belongs to $\varprojlim_NM_N$.

Send the target element to $\sum_{i=1}^de_i\otimes(m_{i,N})_N$.
This is a finite tensor sum and defines a map inverse to $\lambda$:
both composites preserve all basis coordinates at all levels.
For $d=0$ both sides are zero, giving the same conclusion.
\end{proof}

The same conclusion holds when $P$ is a direct summand of a finite
free module, the definition of a finite projective module here.
Write $P\oplus Q\cong S^d$ and apply the theorem to this direct sum.
Tensoring and taking compatible pairs both preserve the two summands.
The resulting isomorphism is the direct sum of the two comparison maps,
so each comparison map is an isomorphism.
The infinite-dimensional factor $L$ in
Proposition~\ref{prop:cc-scalar-extension} shows why flatness alone
cannot replace the finiteness hypothesis.

\section{Localization commutes with quotients}
\label{sec:cc-localization-limit}

Inverting $h$ in a quotient where a power of $h$ is zero forces
every element to vanish, while localizing the limit can leave a
nonzero field. At a single quotient, localization commutes with
forming that quotient by an explicit isomorphism.

Let $S$ be a commutative ring, let $s\in S$, and let $I\subset S$
be an ideal. Write $I[1/s]$ for the ideal of $S[1/s]$ consisting
of fractions with numerator in $I$.
It is an ideal because addition uses a common denominator and a
product with any fraction still has numerator in $I$.

\begin{proposition}\label{prop:cc-localization-quotient}
There is an isomorphism of rings
\[
 (S/I)[1/\bar s]\longrightarrow S[1/s]/I[1/s],\qquad
 \frac{\bar a}{\bar s^m}\longmapsto
       \left[\frac{a}{s^m}\right],
\]
where bars denote classes in $S/I$.
It includes the zero-ring case when a power of $s$ belongs to $I$.
\end{proposition}

\begin{proof}
Replacing $a$ by $a+i$, with $i\in I$, changes the proposed image
by a fraction in $I[1/s]$.
If two fractions agree in $(S/I)[1/\bar s]$, the localization equality
gives $s^\ell(s^na-s^mb)\in I$ for some $\ell$.
Their difference in $S[1/s]$ therefore has a representative with
numerator in $I$, so the proposed images agree.
The map preserves ring operations by the fraction formulas.

In the other direction, the map $a/s^m\mapsto\bar a/\bar s^m$
is well-defined by the same equality rule and kills $I[1/s]$.
It thus induces a map on the quotient. Applying either composite
to a represented fraction returns that fraction, proving the isomorphism.
If $s^N\in I$, then $1=s^N/s^N$ belongs to $I[1/s]$.
Both rings in the asserted isomorphism are consequently zero.
\end{proof}

For $R=k[[h]]$, the system $R/h^NR$, $N\geq1$, uses coefficient
reduction as its transitions.
Proposition~\ref{prop:cc-series-limit} gives
\begin{equation}\label{eq:cc-localization-limit}
 \left(\varprojlim_N R/h^NR\right)[1/h]=K,
 \qquad
 \varprojlim_N\bigl((R/h^NR)[1/\bar h]\bigr)=0.
\end{equation}
Indeed, $\bar h^N=0$ at level $N$, so that localized ring is zero
by Proposition~\ref{prop:cc-localization-quotient}.
The projections from the unlocalized limit induce the comparison
map $K\to0$. Its kernel is all of the nonzero field $K$.
This is the explicit failure of localization to commute with this
inverse limit, even though it commutes with each quotient.

\section{Limits of complexes and compatible primitives}
\label{sec:cc-cohomology}

A cycle can be a boundary at every finite order while having no
single primitive with all those reductions.
This question concerns actual elements of complexes, in addition to
the scalar identities in \eqref{eq:cc-localization-limit}.
An inverse system of cochain complexes $Q_N$ has transitions
$\rho_N:Q_{N+1}\to Q_N$ which preserve degree and satisfy
$d_N\rho_N=\rho_Nd_{N+1}$ as an equality.
Define its limit degree by degree:
\[
 Q^q=\varprojlim_NQ_N^q,\qquad
 d((v_N)_N)=(d_Nv_N)_N.
\]

Commutation with the transitions makes this differential well-defined,
and $d^2=0$ holds at every coordinate.
The underlying graded module is $\bigoplus_{q\in\mathbb Z}Q^q$.
All cohomology here is algebraic cohomology, with the ordinary image
of the differential as its boundary submodule.

The projections of complexes induce the comparison
\[
 \eta^q:H^q\left(\varprojlim_NQ_N\right)
       \longrightarrow\varprojlim_NH^q(Q_N),\qquad
 [(v_N)_N]\longmapsto([v_N])_N.
\]
A compatible cycle has compatible classes. A compatible boundary
maps to zero in every coordinate, proving that $\eta^q$ is well-defined.

Finite free terms with specified compatible bases do reconstruct a
complex before this cohomology question is considered.
Fix integers $a\leq b$. Suppose $R=k[[h]]$ and
$P_N^q=(R/h^NR)^{d_q}$ for $a\leq q\leq b$,
with fixed finite ranks $d_q$, and zero terms outside this interval.
Assume each transition is coordinate reduction in the chosen bases,
and the matrices $D_N^q$ of the differentials reduce compatibly.
Each matrix entry then specifies a unique scalar series in $R$.
Write $D^q$ for the resulting matrix.

Coefficient reduction identifies $\varprojlim_NP_N^q$ with $R^{d_q}$.
Every entry of $D^{q+1}D^q$ reduces to zero at every level, so it
is zero by $\bigcap_Nh^NR=0$.
Thus the limit is the finite free complex with these terms and
differentials. Its reduction modulo $h^N$ recovers exactly $P_N$.
A cohomology comparison remains a separate question.

\begin{example}\label{ex:cc-cohomology-failure}
For $N\geq1$, take
\[
 Q_N=[k[x]\xrightarrow{\ \pi_N\ }k[x]/(x^N)]
\]
in degrees zero and one, with $\pi_N$ the quotient map.
The transition is the identity on degree zero and reduction on
degree one. They commute with $\pi_N$ because both composites take
a polynomial to its class modulo $x^N$.
Both transition maps are surjective.

Each differential is surjective, so
$H^1(Q_N)=0$, while $H^0(Q_N)=x^Nk[x]$.
The induced transition on $H^0$ is the inclusion
$x^{N+1}k[x]\hookrightarrow x^Nk[x]$.
The degree-zero limit is $k[x]$, since compatibility in a constant
system with identity transitions requires a single fixed polynomial.
The degree-one limit is $k[[x]]$, by its compatible coefficients.
The limit differential is the inclusion sending a polynomial to the
same formal series. Consequently
\[
 \varprojlim_NQ_N=[k[x]\hookrightarrow k[[x]]],\qquad
 H^1(\varprojlim_NQ_N)=k[[x]]/k[x]\ne0,
 \qquad\varprojlim_NH^1(Q_N)=0.
\]

The nonzero class of $z(x)=\sum_{j\geq0}x^j$ lies in the kernel of
$\eta^1$. It is nonzero because no polynomial has coefficient one
at every power of $x$.

At level $N$ a primitive is $p_N=\sum_{j=0}^{N-1}x^j$.
These choices satisfy $p_{N+1}-p_N=x^N$ and are not compatible in
degree zero, whose transition is the identity.
Any compatible family of primitives would be a single polynomial
$p$ with $p\equiv z(x)\pmod{x^N}$ for every $N$.
Choose $d\geq0$ such that $p$ has no terms of degree greater than $d$.
At level $N=d+2$, the coefficients of $x^{d+1}$ in $p$ and $z(x)$ differ.
Thus no alternative choice supplies compatible primitives.

This example proves that degreewise surjectivity of the transitions
does not make $\eta^1$ injective.
The complexes $Q_N$ are not acyclic, since their degree-zero
cohomology is nonzero. The example concerns ordinary cohomology:
the image $k[x]$ is dense in $k[[x]]$ for the $x$-adic topology,
so quotienting by the closure of that image would give a different result.
\end{example}

A compatible homotopy supplies primitives at every order simultaneously.
Let $S$ be a commutative ring, let $Q$ be a cochain complex of
$S$-modules, and let $H$ be a graded $S$-module with zero differential.
A contraction of $Q$ onto $H$ consists of degree-zero chain maps
$i:H\to Q$, $p:Q\to H$, and $S$-linear maps
$s^q:Q^q\to Q^{q-1}$ such that
\begin{equation}\label{eq:cc-contraction}
 pi=1_H,\qquad di=0,\qquad pd=0,\qquad
 ds+sd=1_Q-ip.
\end{equation}
For a cycle $v$, this gives $v-ipv=d(sv)$.
If $i(a)=db$, applying $p$ gives $a=0$.
Therefore the map $a\mapsto[i(a)]$ identifies $H^q$ with $H^q(Q)$.

\begin{proposition}\label{prop:cc-compatible-contractions}
Suppose an inverse system $Q_N$ has contractions $(i_N,p_N,s_N)$
onto zero-differential graded modules $H_N$.
Let their transitions be $\rho_N:Q_{N+1}\to Q_N$ and
$\sigma_N:H_{N+1}\to H_N$.
If
\[
 \rho_Ni_{N+1}=i_N\sigma_N,\qquad
 \sigma_Np_{N+1}=p_N\rho_N,\qquad
 \rho_Ns_{N+1}=s_N\rho_N,
\]
then the contractions induce a contraction of $\varprojlim_NQ_N$
onto $\varprojlim_NH_N$, with every limit taken degreewise.
In particular, every cohomology comparison $\eta^q$ is an isomorphism.
\end{proposition}

\begin{proof}
The three compatibility equations ensure that the coordinate maps
$(i_N)$, $(p_N)$, and $(s_N)$ send compatible families to compatible
families of the stated degrees.
Equation~\eqref{eq:cc-contraction} holds at each coordinate and thus
on these limits. Its cycle and boundary argument identifies the
limit cohomology with $(\varprojlim_NH_N)^q$.
At every finite level it also identifies $H_N^q$ with $H^q(Q_N)$.
These identifications commute with transitions by the equation for
$i_N$. Evaluation on a compatible element shows that $\eta^q$ is
exactly the resulting isomorphism between the two descriptions.
\end{proof}

For a cochain complex $Q$, its parameter completion always means
\[
 (Q[[h]])^q=\prod_{n\geq0}h^nQ^q,\qquad
 d\left(\sum_{n\geq0}h^nv_n\right)
       =\sum_{n\geq0}h^ndv_n.
\]
The underlying graded module is the direct sum of these spaces over
$q$. There is no additional product over cochain degrees.
Since $|h|=0$, each coefficient differential has degree one.

\begin{theorem}\label{thm:cc-contraction-completion}
A contraction $(i,p,s)$ of $Q$ onto $H$ extends coefficientwise to
a contraction of $Q[[h]]$ onto $H[[h]]$.
All extended maps are continuous for the degreewise $h$-adic topologies.
Consequently
\[
 H^q(Q[[h]])\cong H^q[[h]]
\]
via $[v]\mapsto p(v)$, where the $H^q$ on the right is the degree-$q$
component of the specified zero-differential module $H$.
\end{theorem}

\begin{proof}
Apply $i,p,s$ to each coefficient, preserving their degrees.
The maps preserve the $h^N$ submodules, so they are continuous.
Every identity in \eqref{eq:cc-contraction} holds at every coefficient
and therefore on the series modules.
For a cycle $v$, the same identity gives $v-ipv=d(sv)$.
For $i(a)=db$, applying $p$ gives $a=0$.
These prove the asserted mutually inverse cohomology maps directly.

The truncations of the maps also satisfy the compatibility conditions
of Proposition~\ref{prop:cc-compatible-contractions}.
\end{proof}

Over a field, such a contraction exists for every cochain complex.
To construct it, write $Z^q=\ker(d:Q^q\to Q^{q+1})$ and
$B^q=\operatorname{im}(d:Q^{q-1}\to Q^q)$.
Lemma~\ref{lem:cc-basis-extension} supplies simultaneous vector-space
complements in every degree:
\[
 Z^q=B^q\oplus H_0^q,\qquad Q^q=Z^q\oplus L^q.
\]
The map $d:L^q\to B^{q+1}$ is injective, since $L^q\cap Z^q=0$.
It is surjective because decomposing a primitive into its $Z^q$ and
$L^q$ parts does not change its differential.

Let $H=H^*(Q)$ with zero differential.
The quotient map identifies $H_0^q$ with $H^q(Q)$, so its inverse
defines $i$. Define $p$ by that quotient on $H_0^q$ and by zero on
$B^q\oplus L^q$.
Define $s$ on $B^q$ by the inverse of $d:L^{q-1}\to B^q$, and
set $s=0$ on $H_0^q\oplus L^q$.
On these three summands, $ds+sd$ equals the identity, zero, and
the identity, respectively. Thus it equals $1-ip$, and the remaining
identities in \eqref{eq:cc-contraction} follow from the definitions.

It follows that, for every complex of $k$-vector spaces with this
degreewise parameter completion, the coefficient map is an isomorphism
\[
 H^q(Q[[h]])\longrightarrow H^q(Q)[[h]],\qquad
 \left[\sum_{n\geq0}h^nv_n\right]\longmapsto
       \sum_{n\geq0}h^n[v_n].
\]
The map itself does not depend on the chosen complements.
The contraction proves its surjectivity by representatives and its
injectivity by coefficientwise primitives.
The varying quotients in Example~\ref{ex:cc-cohomology-failure} do not
have this common coefficient construction, as their incompatible
primitives demonstrate.
